\documentclass[10pt,a4paper]{amsart}
\usepackage[margin=19mm]{geometry}
\usepackage{amsmath,amssymb,mathtools}
\usepackage[scr=rsfs,cal=euler]{mathalfa}
\usepackage{xfrac}
\usepackage[unicode,hidelinks,bookmarks=false]{hyperref}
\newcommand{\ie}{i.e.\ }
\newcommand{\cf}{cf.\ }
\newcommand{\resp}{respectively\ }
\newcommand{\Subsection}{\subsection}
\providecommand{\nobreakdash}{\nobreak\hbox{-}}
\setlength{\emergencystretch}{2em}
\multlinegap0pt
\usepackage{amssymb}
\usepackage{xcolor}
\usepackage{algorithm}\usepackage{algpseudocode}
\usepackage{tikz}
\usepackage[shortlabels]{enumitem}
\newtheorem{theorem}{Theorem}
\title{3-manifold polynomials}
\author{Jos\'e Fr\'ias, Jos\'e Carlos G\'omez-Larra\~naga, Jos\'e Luis Le\'on-Medina, Fabiola Manjarrez-Guti\'errez}
\date{Source: arXiv:2510.06651v2 (CC BY 4.0)}
\begin{document}
\maketitle

\noindent\emph{Source and licence.} 3-manifold polynomials. Version arXiv:2510.06651v2 (CC BY 4.0), distributed by arXiv under the Creative Commons Attribution 4.0 International licence, http://creativecommons.org/licenses/by/4.0/, as recorded in the arXiv licence metadata for this exact version and shown on the abstract page https://arxiv.org/abs/2510.06651v2. Full licence notice: \texttt{licenses/arxiv-2510.06651-CC-BY-4.0.txt}.\par\smallskip

\noindent\textit{Editorial scope.} Counted unit: the article's theorem thm:square-shape (source0.tex lines 634--647). The article's complete original proof of it is delivered with it (source0.tex lines 649--682, one \texttt{proof} environment of 34 lines). No mathematics is written, repaired or re-derived: the statement and the proof are the article's own lines, in the article's own notation, and no retained line is substituted or re-flowed.  Retained uncounted as the source-local context the proof needs: lines 613--619.  Nothing was omitted from any retained span, so the package carries no omission record.  No retained line contains a citation, so the wrapper carries no bibliography and no bibliography entry.  No retained line contains a figure environment, an image-inclusion command or drawing code, so no image is delivered and the package declares no asset dependency.  Editorial formatting only, in addition: an AMS \texttt{amsart} preamble that restates the article's own declarations and macros used by the retained lines, namely the environments \texttt{theorem} on separate counters, together with the standard packages those lines need.  The retained environments print in this document as Theorem 1 in order of appearance, whereas the article prints the same items with the numbering of its own full document.  The complete native source of the article as distributed in the arXiv e-print for this exact version is bundled unchanged as \texttt{sources/186623/source0.tex}.
\begin{theorem}\label{thm:product-and-parity}
Let $G=C_p(\pm 1,\pm q)$ with $\gcd(p,q)=1$, $p\ge3$. Then
\[
\tau(p,q)=\frac{1}{p}\prod_{j=1}^{p-1}\Bigl(4-2\cos\tfrac{2\pi j}{p}-2\cos\tfrac{2\pi q j}{p}\Bigr),
\]
and $\tau(p,q)$ depends only on the orbit $\{\pm q^{\pm1}\}\subset(\mathbb{Z}/p\mathbb{Z})^\times$.
\end{theorem}

\begin{theorem}\label{thm:square-shape}
Let $G=C_p(\pm 1,\pm q)$ with $\gcd(p,q)=1$, $p\ge3$. Then:
\begin{itemize}
\item[\emph{(a)}] If $p$ is odd, there exists an integer $A$ (depending on $p$, $q$) such that
\[
\tau(p,q)=p\,A^2.
\]
\item[\emph{(b)}] If $p$ is even, there exists an integer $B$ (depending on $p$, $q$) such that
\[
\tau(p,q)=\frac{\lambda_{p/2}}{p}\,B^2,
\qquad \lambda_{p/2}=6-2(-1)^q\in\{4,8\}.
\]
\end{itemize}
\end{theorem}

\begin{proof}
By Theorem~\ref{thm:product-and-parity},
\[
\tau(p,q)=\frac{1}{p}\prod_{j=1}^{p-1}\lambda_j,
\qquad 
\lambda_j=4-2\cos\frac{2\pi j}{p}-2\cos\frac{2\pi q j}{p}.
\]

\begin{itemize}
    \item[(a)] Suppose $p$ is odd. Hence, the product pairs up symmetrically:
\[
\prod_{j=1}^{p-1}\lambda_j=\left(\prod_{j=1}^{(p-1)/2}\lambda_j\right)^2,
\]
hence
\[
\tau(p,q)=\frac{1}{p}\left(\prod_{j=1}^{(p-1)/2}\lambda_j\right)^2.
\]
As $\tau(p,q)\in\mathbb{Z}$, the product must be divisible by $p$, say $\prod_{j=1}^{(p-1)/2}\lambda_j=p\,A$, giving $\tau(p,q)=p\,A^2$.

\item[(b)] Suppose $p$ is even. All eigenvalues pair except the unpaired one at $j=p/2$:
\[
\prod_{j=1}^{p-1}\lambda_j=\lambda_{p/2}\left(\prod_{j=1}^{(p/2)-1}\lambda_j\right)^2.
\]
Since,
\[
\lambda_{p/2}=4-2\cos\pi-2\cos(q\pi)=6-2(-1)^q\in\{4,8\}.
\]
We get,
\[
\tau(p,q)=\frac{\lambda_{p/2}}{p}\,B^2,
\qquad B=\prod_{j=1}^{(p/2)-1}\lambda_j \in\mathbb{Z}.
\]  
\end{itemize}
\end{proof}

\end{document}
